% Configurazione
\documentclass{article}

\usepackage{titling} % Required for inserting the subtitle
\usepackage{graphicx} % Required for inserting images
\usepackage{tabularx} % Per l'ambiente tabularx (tabelle)
\usepackage{calc} % Sempre per le tabelle
\usepackage[hidelinks]{hyperref} % Per i collegamenti ipertestuali, ad esempio sulla table of contents
\usepackage[italian]{babel} % Per la lingua italiana nelle scritte automatiche
\usepackage{xcolor} % Per colorare il testo
\usepackage{colortbl} % Per colorare le celle delle tabelle
\usepackage{lipsum} % Per generare lorem ipsum
\usepackage[normalem]{ulem} % Per sottolineare il testo
\usepackage{array} % Per la visualizzazione fluttuante di array di domande e risposte
\usepackage{ragged2e} % Pacchetto necessario per \justifying che giustifica il testo di tabelle

\newcommand{\ulhref}[2]{\href{#1}{\uline{#2}}} % Nuovo comando per sottolineare i link
\newcommand{\ulref}[1]{\uline{\ref{#1}}} % Nuovo comando per sottolineare i collegamenti a immagini

\graphicspath{ {immagini/} {../../shared/immagini/} }


% Struttura
\begin{document}

\begin{minipage}{0.4\textwidth}
    \includegraphics[width=0.6\textwidth]{Logo Universita di Padova.png}
\end{minipage}
\begin{minipage}{0.55\textwidth}
    \textcolor{red}{\textbf{Università degli Studi di Padova}} \\
    \textcolor{red}{Laurea: Informatica} \\
    \textcolor{red}{Corso: Ingegneria del Software} \\
    \textcolor{red}{Anno Accademico: 2024/2025}
\end{minipage}

\begin{minipage}{0.4\textwidth}
    \includegraphics[width=0.6\textwidth]{Logo SWEg.png}
\end{minipage}
\begin{minipage}{0.55\textwidth}
    \textbf{Gruppo: SWEg Labs} \\
    Email: \textsf{gruppo.sweg@gmail.com}
\end{minipage}

\vspace{2cm}

{
\centering
\Huge\bfseries Lettera di Presentazione\par
}

\vspace{2cm}
\section{Informazioni generali}

\begin{itemize}
    \item \textbf{Tipo di riunione:} esterna con proponenti
    \item \textbf{Luogo:} meeting zoom
    \item \textbf{Data:} 22/10/2024
    \item \textbf{Ora inizio:} 17:00
    \item \textbf{Ora fine:} 17:40
    \item \textbf{Responsabile:} Stefani Riccardo
    \item \textbf{Notaio:} Verzotto Davide
    \item \textbf{Partecipanti:}
    \begin{itemize}
        \renewcommand{\labelitemii}{--}
        \item Bolognini Federica
        \item Fantinato Michael
        \item Loat Giacomo
        \item Righetto Filippo
        \item Stefani Riccardo
        \item Verzotto Davide
        \item Carlesso Gianluca (\emph{Ergon Informatica Srl})
        \item Tieppo Anna (\emph{Ergon Informatica Srl})
    \end{itemize}
\end{itemize}

\newpage
\section*{Registro delle modifiche}

\begin{table}[h]
    \centering
    \begin{tabular}{|l|l|l|p{5cm}|}
        \hline
        \rowcolor[gray]{0.9}
        \textbf{Vers.} & \textbf{Data} & \textbf{Autore} & \textbf{Descrizione}\\
        \hline
        1.0 & ... & ... & Approvazione del documento\\
        \hline
        0.9 & ... & ... & Verifica del documento\\
        \hline
        0.7 & ... & ... & ...\\
        \hline
        0.6 & ... & ... & ...\\
        \hline
        0.5 & ... & ... & ...\\
        \hline
        0.2 & 25-10-2024 & Riccardo Stefani & Scrittura della motivazione della scelta e della conclusione\\
        \hline
        0.1 & 25-10-2024 & Riccardo Stefani & Creazione del documento\\
        \hline
    \end{tabular}
\end{table}

\newpage
\tableofcontents
\newpage

\section{Introduzione}

In questo documento viene presentata l'analisi dei capitolati d'appalto che abbiamo svolto al fine di valutarne la fattibilità e scegliere il più adatto 
per il gruppo di lavoro \emph{SWEg Labs}. Verranno esposte in maniera chiara le motivazioni alla base della scelta di candidarci per il capitolato 
\emph{BuddyBot}, e altrettanto verranno analizzate le altre due proposte che avevamo fortemente considerato, \emph{LLM: Assistente digitale} e 
\emph{Vimar GENIALE}. Verranno altresì elencati i pro e i contro di ciascuno dei capitolati rimanenti per argomentare la decisione di averli scartati.
\section{Valutazione del Capitolato selezionato}

\subsection{C9: BuddyBot}

\subsubsection{Descrizione}
Il capitolato richiede di sviluppare un LLM capace di rispondere
a diversi tipi di richieste relative ai dati interni all'azienda, 
rendendo inoltre possibile per gli utenti ricevere informazioni riguardanti 
i contenuti presenti su piattaforme esterne.

\subsubsection{Dominio}
\paragraph{Dominio tecnologico:}
Viene richiesto l'utilizzo di un LLM open-source o di API di modelli 
a pagamento (come ChatGPT). Il modello deve poter comunicare tramite le API di Jira, 
Confluence e GitHub. Una volta recuperate le informazioni necessarie per la risposta,
deve elaborarle e formulare una risposta finale. Tutte le chat devono essere
salvate su un database e devono poter essere recuperate. L'azienda non richiede
l'uso di una tecnologia specifica, ma fornisce alcuni suggerimenti:
\begin{itemize}
    \item \textit{\textbf{OpenAI}}: Utilizzo tramite API per le funzionalità NLP, ovvero di comprensione e generazione del testo.
    \item \textit{\textbf{LangChain}}: Progetto open-source che permette l'integrazione di modelli AI senza richiedere una conoscenza dettagliata del loro funzionamento interno.
    \item \textit{\textbf{Angular}}: Framework front-end per la costruzione di applicazioni web.
    \item \textit{\textbf{Node/NestJS}}: Framework per lo sviluppo di applicazioni server-side orientato ai microservizi e alle API RESTful.
    \item \textit{\textbf{Spring Boot}}: Framework Java che offre strumenti per la creazione di applicazioni standalone, con supporto integrato per database, sicurezza e gestione delle dipendenze.
\end{itemize}


\paragraph{Dominio applicativo:}
BuddyBot vuole diventare un assistente a 360 gradi per qualsiasi membro 
dell'azienda. Si propone di essere facilmente accessibile tramite
un'interfaccia web, dove l'utente può formulare domande e recuperare lo storico delle vecchie chat.
Il bot deve poter rispondere a domande riguardanti file presenti in Jira, GitHub e Confluence, 
riuscendo anche a incrociare dati provenienti da diverse fonti. Il compito principale di 
BuddyBot sarà quindi aiutare i membri dell'azienda a trovare e combinare i dati in 
modo più rapido e preciso, fornendo risposte in linguaggio naturale.

\subsubsection{Motivazione della scelta}
La scelta finale è ricaduta su questo capitolato per diversi motivi:
\begin{itemize}
    \item \textit{\textbf{Disponibilità dell'azienda}}: L'azienda si è dimostrata da subito molto disponibile; le risposte alle mail sono arrivate in brevissimo tempo e si è subito manifestata la volontà di organizzare un incontro. Inoltre, si sono proposti di fornire anche supporto tecnico durante lo sviluppo, qualora ce ne fosse bisogno.
    \item \textit{\textbf{Interesse}}: Il prodotto finale ha suscitato grande interesse nel gruppo, soprattutto per le tecnologie con cui prevediamo di lavorare per realizzarlo. In particolare, ci è sembrato il capitolato con il miglior utilizzo dell'IA.
    \item \textit{\textbf{Fattibilità del progetto}}: Confrontando le richieste tra i vari capitolati, BuddyBot ci è sembrato il miglior compromesso tra difficoltà e interesse.
    \item \textit{\textbf{Libertà tecnologica}}: La libertà lasciata nella scelta delle tecnologie ci è sembrata molto vantaggiosa, permettendoci di utilizzare tecnologie con le quali alcuni membri del gruppo hanno già esperienza.
\end{itemize}

\subsubsection{Conclusioni}
Visti i motivi sopra citati e non trovando in nessun altro capitolato una combinazione equivalente, 
il gruppo ha deciso di scegliere BuddyBot come progetto per cui candidarsi.

\section{Valutazione dei Capitolati papabili}


\subsection{C7: LLM: Assistente digitale}

\subsubsection{Descrizione}

\subsubsection{Dominio}

\subsubsection{Criticità riscontrate}

\subsubsection{Conclusioni}


\subsection{C2: Vimar GENIALE}

\subsubsection{Descrizione}

\subsubsection{Dominio}

\subsubsection{Criticità riscontrate}

\subsubsection{Conclusioni}
\section{Valutazione degli altri Capitolati}


\subsection{C1: Artificial QI}

\subsubsection{Pro}

\begin{itemize}
    \item \textbf{Innovazione e utilità}: il gruppo trova interessante cercare di trovare una soluzione per il problema molto attuale della valutazione delle abilità degli LLM.
    \item \textbf{Punti di partenza}: il proponente Zucchetti è disponibile a fornire un programma esterno che espone delle API da testare, che è un'opportunità per avere una parte del sistema già pronta.
    \item \textbf{Molti requisiti opzionali}: sono consigliati vari spunti per estendere l'applicazione nel caso in cui si riesca a concludere in tempi brevi i requisiti obbligatori.
\end{itemize}

\subsubsection{Contro}

\begin{itemize}
    \item \textbf{Complessità del problema}: il capitolato tratta un problema aperto nel mondo dell'informatica, il che costituisce una complessità di ricerca esplorativa che va oltre lo sviluppo dell'applicativo e che non è calcolabile a priori.
    \item \textbf{Genericità degli strumenti}: vengono presentate molte poche tecnologie consigliate rispetto ad altri capitolati, e i suggerimenti rimangono molto generici a riguardo.
    \item \textbf{Dominio incognito e dataset assente}: non è chiaro quale sia il dominio di conoscenza che è da testare nell'LLM di turno, infatti non viene mai specificato che tipo di domande e risposte devono essere presenti nel dataset. Sembra inoltre che sia compito del gruppo scrivere le domande e risposte che facciano da training set, il che costituisce un onere non da poco considerata l'enorme quantità di dati necessari per allenare il sistema.
\end{itemize}


\subsection{C3:	Automatizzare le routine digitali tramite l’intelligenza generativa}

\subsubsection{Pro}

\begin{itemize}
    \item \textbf{Automazione delle Routine}: È interessante che il progetto punti ad automatizzare compiti ripetitivi come l’organizzazione delle e-mail e la gestione dei calendari, migliorando così l’efficienza e l'organizzazione aziendale.
    \item \textbf{Estendibilità}: L’applicativo è concepito in modo modulare, permettendo di aggiungere facilmente nuove automazioni e blocchi funzionali.
    \item \textbf{Supporto e Formazione}: Var Group fornisce supporto continuo per le tecnologie utilizzate, formazione e incontri periodici per garantire la qualità del progetto e per gestire l'evoluzione del software.
\end{itemize}

\subsubsection{Contro}

\begin{itemize}
    \item \textbf{Incremento dei tempi di sviluppo}: Potremmo dover affrontare ritardi dovuti alla necessità di studiare nuove e diverse tecnologie tra cui Python, Swift, NodeJS, MongoDB e Typescript. Questo potrebbe impattare la consegna nei tempi previsti, con conseguente aumento dei costi di progetto.
    \item \textbf{Scarso interesse}: il gruppo ha preferito concentrarsi su altri progetti con focus più accattivanti.
\end{itemize}


\subsection{C4:	NearYou - Smart custom advertising platform}

\subsubsection{Pro}

\begin{itemize}
    \item \textbf{Innovatività}: il progetto utilizza tecnologie avanzate come intelligenza artificiale generativa, data stream processing e targeting geografico in tempo reale. Questo lo rende un progetto ideale per acquisire esperienza con tecnologie molto richieste nel mercato attuale.
    \item \textbf{Flessibilità tecnologica}: il progetto lascia ampia libertà di scelta sugli strumenti e tecnologie da adottare.
    \item \textbf{Licenza}: la licenza del software sarà interamente nostra, potendo esporre il progetto nel curriculum.
\end{itemize}

\subsubsection{Contro}

\begin{itemize}
    \item \textbf{Complessità elevata}: il progetto richiede la gestione di argomenti forse troppo complessi data la scarsa esperienza dei membri del gruppo negli ambiti richiesti.
    \item \textbf{Rischio di dispersione}: la possibilità di scegliere le tecnologie da utilizzare potrebbe essere un'arma a doppio taglio, portando a dispersione o a scelte tecniche inadeguate, soprattutto se non si è esperti in tutti gli ambiti.
    \item \textbf{Scarso interesse}: il gruppo ha preferito concentrarsi su altri progetti con focus più accattivanti.
\end{itemize}


\subsection{C5:	3Dataviz}

\subsubsection{Pro}

\begin{itemize}
    \item \textbf{Gestione dei dati}: Il progetto è molto utile per quanto riguarda l’apprendimento, la gestione e la visualizzazione di un grande numero di dati in maniera chiara e interattiva.
    \item \textbf{Ampia libertà progettuale}: Il progetto offre grande libertà dal punto di vista della realizzazione dei vari grafici e delle funzionalità richieste. Questa flessibilità consente di personalizzare il sito in base alle esigenze e alle competenze del team.
\end{itemize}

\subsubsection{Contro}

\begin{itemize}
    \item \textbf{Richieste competenze avanzate}: La creazione e gestione di un sistema di visualizzazione 3D richiede delle competenze tecniche molto avanzate nel campo della grafica tridimensionale.
    \item \textbf{Studio di tecnologie nuove}: Il progetto richiede lo studio di tecnologie mai utilizzate prima, rendendo necessario un approfondimento e uno studio individuale per acquisire le competenze necessarie.
    \item \textbf{Scarso interesse}: il gruppo ha preferito concentrarsi su altri progetti con focus più accattivanti.
\end{itemize}


\subsection{C6: Sistema di gestione di un magazzino distribuito}

\subsubsection{Pro}

\begin{itemize}
    \item \textbf{Machine Learning}: Si richiede l'utilizzo del machine Learning per prevedere la richiesta futura di un certo materiale, materia che può essere molto interessante su cui lavorare.
    \item \textbf{Utilità}: Il problema della gestione del magazzino c’e’ da molti anni e colpisce molte aziende, imparare a creare software innovativi per la gestione di esso può essere molto utile e formativo.
\end{itemize}

\subsubsection{Contro}

\begin{itemize}
    \item \textbf{Complessità dei microservizi}: L'architettura basata su microservizi introduce maggiore complessità nella gestione del sistema e nell'orchestrazione, richiedendo più competenze e lavoro..
    \item \textbf{Sicurezza}: Richiesta di sicurezza avanzata tra il sistema centrale e i sistemi locali, non richiesto da altri capitolati. Questa richiesta aumenta la complessità del sistema, rendendo anche più difficile quantificare le ore di allenamento necessarie.
    \item \textbf{Requisiti di test e validazione}: La necessità di effettuare test estensivi e validazioni di algoritmi predittivi e di sicurezza, implica un elevato sforzo continuo.
\end{itemize}


\subsection{C8: Requirement Tracker - Plug-in VS Code}

\subsubsection{Pro}

\begin{itemize}
    \item \textbf{Supervisione continua}: disponibilità dell’azienda ad essere presente sia in presenza che in remoto, potendo così assistere il gruppo e guidarlo nella realizzazione.
    \item \textbf{Modelli di AI pre-addestrati}: è possibile utilizzare modelli di AI già esistenti e sfruttarli tramite API. Il team potrebbe quindi iniziare a lavorare senza dover programmare da zero un modello.
\end{itemize}

\subsubsection{Contro}

\begin{itemize}
    \item \textbf{Necessario studio di tecnologie sconosciute}: è necessario implementare un’architettura di tipo modulare, che il gruppo non ha ancora usato approfonditamente, dunque nonostante l’occasione, il suo apprendimento rallenterebbe il lavoro per il tempo necessario.
    \item \textbf{Vincoli stringenti}: i vincoli obbligatori imposti sono in quantità tale da presentare un ridotto spazio di manovra o di inventiva per il team.
    \item \textbf{Proof of Concept}: la necessità di sviluppare un P.o.C. già a metà progetto, può richiedere uno sforzo elevato, soprattutto essendo previste numerose tecnologie nuove da apprendere.
\end{itemize}

\end{document}
